% Gauge Structure and Effective Geometry from Admissible Non-Injective Projection
% Cosmochrony companion paper M4

\documentclass[11pt,a4paper]{article}

\usepackage{amsmath,amssymb,amsthm}
\usepackage{enumerate}
\usepackage{hyperref}
\usepackage{doi}

\newtheorem{theorem}{Theorem}[section]
\newtheorem{lemma}[theorem]{Lemma}
\newtheorem{proposition}[theorem]{Proposition}
\newtheorem{corollary}[theorem]{Corollary}
\newtheorem{definition}[theorem]{Definition}
\newtheorem{remark}[theorem]{Remark}

\title{Gauge Structure and Effective Geometry\\
from Admissible Non-Injective Projection}

\author{J\'er\^ome Beau\\
\small Independent Researcher\\
\small \texttt{jerome.beau@cosmochrony.org}}

\date{2026}

\begin{document}

  \maketitle

  \begin{abstract}
    Starting from the axiomatic foundation of admissible non-injective
    transitions~\cite{FoundationM}, we derive the macroscopic structures of gauge
    theory and effective spacetime geometry without introducing them as postulates.
    The derivation proceeds in three steps.
    First, gauge charges are identified as invariants of equivalence classes in
    the projection fibre: two configurations are gauge-equivalent when they project
    to the same observable state.
    The group of fibre-preserving transformations constitutes the gauge group,
    with $\mathrm{U}(1)$, $\mathrm{SU}(2)$, and $\mathrm{SU}(3)$ identified as
    natural candidates from the admissibility structure.
    Second, in the continuum limit $q \to \infty$ established in~\cite{FoundationM},
    the discrete admissibility filter $\Pi_n$ converges to an effective continuum
    operator $\mathcal{L}_\Pi$ whose principal symbol defines an effective metric
    tensor $g^{\mu\nu}(x) \propto A^{\mu\nu}(x)$.
    Third, the bounded Born--Infeld admissibility constraint (A3--A4) restricts
    the admissible forms of $A^{\mu\nu}(x)$: in homogeneous regimes, hyperbolicity
    and spectral stability uniquely select Minkowski geometry; in the presence of a
    localized stationary obstruction, flux conservation forces the Schwarzschild
    geometry as the unique stationary exterior solution.
    The gauge connection arises as the correction required to compare admissible
    fibres over neighbouring points, and Yang--Mills dynamics follow from minimising
    the curvature of this connection.
    No dynamical substrate is required at any step; the derivations rest solely on
    Axioms~A1--A4 of~\cite{FoundationM} and the spectral results of the O-series.
  \end{abstract}

  \section{Introduction}
    \label{sec:intro}

    \subsection*{Position in the programme}

      The companion papers~\cite{FoundationM,Q1,Q2} establish the quantum-mechanical
      structure of the Cosmochrony framework: the proto-state, phase coherence, the
      Born rule, the singlet correlator, and the Tsirelson bound, all derived from the
      admissibility structure of non-injective projection.
      These results concern the \emph{internal} structure of the admissible fibre at
      a single transition step.

      The present paper addresses the \emph{macroscopic} consequences of allowing the
      projection structure to vary from point to point in the effective continuum limit.
      Two questions drive the paper:
      \begin{enumerate}[(i)]
\item What are the conserved invariants of the projection fibre, and what
symmetry group do they generate?
\item What effective geometry is selected when the admissibility constraints
of~\cite{FoundationM} are imposed on a spatially varying projection?
      \end{enumerate}

      The answers are: gauge charges and gauge groups (§\ref{sec:charges}), and
      Lorentzian geometry with Schwarzschild as the exterior stationary solution
      (§\ref{sec:geometry}).

    \subsection*{What is and is not derived in this paper}

      This paper derives the \emph{structure} of gauge theory and effective geometry
      from the admissibility framework.
      It does not derive specific coupling constants, the precise matter content of
      the Standard Model, or the quantitative value of Newton's constant.
      These remain open and are identified as directions for future work in
      Section~\ref{sec:open}.

      The derivations are structural: they identify which gauge groups and which
      effective geometries are \emph{admissible} given the constraints, not which
      specific values their parameters take.

    \subsection*{Organisation}

      Section~\ref{sec:continuum} establishes the continuum limit of the discrete
      projection $\Pi_n$ and the emergence of an effective differential operator.
      Section~\ref{sec:charges} derives gauge charges as fibre invariants and
      identifies the gauge group.
      Section~\ref{sec:connection} derives the gauge connection from the requirement
      of consistent transport between neighbouring fibres.
      Section~\ref{sec:ym} derives Yang--Mills dynamics from the minimisation of
      connection curvature.
      Section~\ref{sec:geometry} derives the effective metric from the principal
      symbol of $\mathcal{L}_\Pi$ and establishes the emergence of Minkowski and
      Schwarzschild geometries.
      Section~\ref{sec:open} identifies the open directions.

  \section{Continuum Limit of the Admissibility Projection}
    \label{sec:continuum}

    \subsection{From discrete $\Pi_n$ to a continuum operator}

      The discrete projection $\Pi_n$ acts on the orbit span
      $\Omega_n^{(c)} \subset \mathbb{C}^q$, filtering admissible directions via the
      Born--Infeld bound $A_n \leq c_\chi/\sqrt{\lambda_n}$ (Foundation M,
      §\ref{subsec:five-levels}).
      As established in~\cite{FoundationM} Section~3.5, the large-$q$ limit sends:
      \begin{align}
        G_q &\;\longrightarrow\; \mathrm{Heis}_3(\mathbb{R}),\\
        \mathbb{C}^q &\;\longrightarrow\; L^2(\mathbb{R}),\\
        \text{BFS shells } S_n &\;\longrightarrow\; \text{continuous foliation},\\
        \Pi_n \text{ (discrete)} &\;\longrightarrow\; \mathcal{L}_\Pi \text{ (continuum operator)}.
      \end{align}

      The continuum limit of relational Laplacians on graphs of the type
      $\mathrm{Cay}(G_q, S_q)$ is governed by standard spectral convergence
      results~\cite{BornInfeld}: under mild regularity conditions, the graph Laplacian
      converges in the strong-resolvent sense to a second-order differential operator
      \begin{equation}
        \label{eq:continuum-op}
        \mathcal{L}_\Pi \;=\; \nabla_\mu\bigl(A^{\mu\nu}(x)\,\nabla_\nu\bigr),
      \end{equation}
      where $A^{\mu\nu}(x)$ is a symmetric tensor encoding the local connectivity
      and stiffness of the admissible projection.

      \begin{remark}[Status of the continuum limit]
        The convergence $\Pi_n \to \mathcal{L}_\Pi$ is established qualitatively
        in~\cite{FoundationM} and in detail in~\cite{BornInfeld} for the class of
        weighted relational Laplacians.
        A fully rigorous proof at the level of the Heisenberg Cayley graph
        $\mathrm{Cay}(G_q, S_q)$ as $q \to \infty$ remains open.
        The present section uses the continuum limit as a working hypothesis,
        stated honestly as such.
      \end{remark}

    \subsection{The operator-metric correspondence}

      In admissible continuum regimes, the operator $\mathcal{L}_\Pi$ has a
      well-defined principal symbol:
      \begin{equation}
        \label{eq:principal-symbol}
        \sigma_2(\mathcal{L}_\Pi)(x, k)
        \;=\; A^{\mu\nu}(x)\,k_\mu k_\nu.
      \end{equation}
      This symbol defines an effective metric tensor~\cite{BornInfeld}:
      \begin{equation}
        \label{eq:metric-from-symbol}
        g^{\mu\nu}(x) \;\propto\; A^{\mu\nu}(x).
      \end{equation}
      The metric is not an independent dynamical variable: it is a derived object
      encoding how admissible modes propagate in the continuum limit.
      This is the operator-metric correspondence on which all subsequent geometric
      results rest.

  \section{Gauge Charges as Fibre Invariants}
    \label{sec:charges}

    \subsection{Charges from non-injectivity}

      The non-injectivity of $\Pi_n$ (Axiom~A2 of~\cite{FoundationM}) means that
      multiple configurations in $\Omega_n^{(c)}$ project to the same observable
      state $O_n$.
      These configurations are physically indistinguishable at the observable level
      but may carry distinct internal structure.

      \begin{definition}[Gauge charge]
        \label{def:gauge-charge}
        A \emph{gauge charge} is a functional
        \begin{equation}
          Q : \Omega_n^{(c)} \;\longrightarrow\; \mathbb{R}
        \end{equation}
        that is constant on projection fibres:
        \begin{equation}
          \Pi_n(\omega_1) = \Pi_n(\omega_2)
          \;\Longrightarrow\;
          Q(\omega_1) = Q(\omega_2).
        \end{equation}
      \end{definition}

      A gauge charge is therefore an observable that cannot be resolved by $\Pi_n$
      itself — it lives in the kernel structure of $\Pi_n$ — but is well-defined on
      the fibre.

      \begin{definition}[Gauge group]
        \label{def:gauge-group}
        The \emph{gauge group} $\mathcal{G}$ is the group of transformations
        $\omega \mapsto g \cdot \omega$ that preserve the projection:
        \begin{equation}
          \mathcal{G}
          \;:=\;
          \bigl\{g \in \mathrm{Aut}(\Omega_n^{(c)}) \mid \Pi_n(g \cdot \omega) = \Pi_n(\omega)
          \text{ for all } \omega \bigr\}.
        \end{equation}
      \end{definition}

      $\mathcal{G}$ consists precisely of the fibre-preserving symmetries: it acts
      within $\Pi_n^{-1}(O_n)$ without affecting the observable output.

      \begin{remark}[Gauge invariance as structural necessity]
        The gauge group is not introduced by hand.
        It is forced by the structure of $\Pi_n$: as soon as $\Pi_n$ is non-injective,
        there exist non-trivial fibre-preserving transformations, and hence a
        non-trivial gauge group.
        Gauge invariance is a structural consequence of Axiom~A2, not an additional
        postulate.
      \end{remark}

    \subsection{Identification of the gauge groups}

      Three natural candidates emerge from the fibre structure established in
      Q1~\cite{Q1} and Q2~\cite{Q2}:

      \begin{itemize}
        \item \emph{$\mathrm{U}(1)$}: global phase rotations $\omega \mapsto e^{i\theta}\omega$
        within the complex fibre $\mathbb{C}^q$.
        These preserve $|\omega|$ and hence the observable norm, while acting
        non-trivially on the phase — which is in the kernel of $\Pi_n$ by the
        phase invariance of $\sigma_c(n)$ (O12).
        \item \emph{$\mathrm{SU}(2)$}: the stable fixed point of the admissibility flow
        established in~\cite{Q2} Theorem~8.1.
        Its fibre-preserving action on $\mathbb{C}^2$ is the spin-$\tfrac{1}{2}$
        structure derived in~\cite{Q1}.
        \item \emph{$\mathrm{SU}(3)$}: a structurally motivated candidate for a
        higher-dimensional internal symmetry of the fibre, corresponding to
        multiple co-admissible internal directions beyond $\mathrm{SU}(2)$.
        Its identification from the O-series admissibility data is an open direction
        (Section~\ref{sec:open}).
      \end{itemize}

      \begin{remark}[Status of SU(3)]
        The identification of $\mathrm{SU}(3)$ as a gauge group is a structural
        hypothesis, not a derived result.
        Establishing it would require finding, within the O-series spectral data,
        a co-admissibility pattern for a group with $\mathrm{SU}(3)$ structure
        analogous to the $\lambda_{1/2} = \lambda_{3/2}$ co-admissibility of~\cite{Q2}.
        This is identified as a medium-term open direction.
      \end{remark}

  \section{The Gauge Connection}
    \label{sec:connection}

    \subsection{The problem of comparing neighbouring fibres}

      In the continuum limit, the admissible fibre $F_n$ becomes a smooth fibre
      bundle over the effective spacetime manifold.
      For each point $x$ in this manifold, the fibre $F_x = \Pi^{-1}(x)$ carries
      the internal gauge structure of Section~\ref{sec:charges}.

      When the projection $\Pi(x)$ varies from point to point, the fibres
      $F_x$ and $F_{x+dx}$ are not canonically identified.
      A state $\psi(x)$ transforming under $\mathcal{G}$ acquires an additional
      term under differentiation:
      \begin{equation}
        \partial_\mu(g(x)\,\psi(x))
        \;=\; g(x)\,\partial_\mu\psi(x) + (\partial_\mu g(x))\,\psi(x),
      \end{equation}
      which is not covariant under local gauge transformations $g(x) \in \mathcal{G}$.

    \subsection{Gauge connection from fibre transport}

      \begin{definition}[Gauge connection]
        \label{def:connection}
        A \emph{gauge connection} $A_\mu$ is the correction required to define a
        covariant derivative
        \begin{equation}
          \label{eq:covariant-derivative}
          D_\mu \;:=\; \partial_\mu + A_\mu
        \end{equation}
        such that $D_\mu\psi$ transforms homogeneously under local gauge
        transformations.
      \end{definition}

      \begin{proposition}[Admissibility interpretation of $A_\mu$]
        \label{prop:connection-admissibility}
        In the projection framework, $A_\mu(x)$ encodes the local variation of the
        admissible fibre structure: it is the infinitesimal correction needed to
        transport an admissible configuration from the fibre at $x$ to the fibre
        at $x + dx$ consistently with the admissibility constraints.
      \end{proposition}

      \begin{proof}[Structural argument]
        The admissibility constraints (A1--A3 of~\cite{FoundationM}) define, at each
        point $x$, a fibre $F_x$ of structurally compatible configurations.
        When $\Pi(x)$ varies, the fibres $F_x$ and $F_{x+dx}$ are related by the
        differential of the projection map.
        Consistency of the admissibility condition across neighbouring fibres —
        i.e. that a configuration admissible at $x$ maps to an admissible configuration
        at $x+dx$ — requires a linear correction of the form $A_\mu(x)\,dx^\mu$.
        This correction is precisely the gauge connection.
      \end{proof}

      \begin{remark}[Relation to the holonomy argument of the companion paper]
        The white paper~\cite{Cosmochrony} establishes that non-injective admissible
        projections generically have non-trivial holonomy, which is the minimal
        structural condition for a gauge symmetry to arise as a descriptive invariance.
        Proposition~\ref{prop:connection-admissibility} is the local infinitesimal
        version of this global statement.
      \end{remark}

  \section{Yang--Mills Dynamics}
    \label{sec:ym}

    \subsection{Field strength as curvature of the connection}

      \begin{definition}[Field strength]
        \label{def:field-strength}
        The \emph{field strength} associated with the gauge connection $A_\mu$ is
        the curvature of the covariant derivative:
        \begin{equation}
          \label{eq:field-strength}
          F_{\mu\nu} \;:=\; [D_\mu, D_\nu]
          \;=\; \partial_\mu A_\nu - \partial_\nu A_\mu + [A_\mu, A_\nu].
        \end{equation}
      \end{definition}

      $F_{\mu\nu}$ measures the failure of commutativity of parallel transport around
      an infinitesimal loop — equivalently, the non-commutativity of the admissibility
      constraints in two neighbouring directions.

    \subsection{Yang--Mills action from admissibility minimisation}

      \begin{proposition}[Yang--Mills action from coherence minimisation]
        \label{prop:ym-action}
        The natural dynamical principle for the gauge connection is minimisation
        of the deformation of the projection structure, leading to the action
        \begin{equation}
          \label{eq:ym-action}
          S_{\mathrm{YM}}[A]
          \;=\; \int \mathrm{Tr}\bigl(F_{\mu\nu} F^{\mu\nu}\bigr)\,\sqrt{-g}\,d^4x.
        \end{equation}
      \end{proposition}

      \begin{proof}[Structural argument]
        By Proposition~\ref{prop:connection-admissibility}, $A_\mu$ encodes the local
        variation of the fibre structure.
        The curvature $F_{\mu\nu}$ measures the incompatibility of this variation
        around closed loops: a flat connection ($F_{\mu\nu} = 0$) corresponds to a
        globally consistent fibre structure, while non-zero curvature signals
        inconsistency.
        The admissibility principle (A3 of~\cite{FoundationM}) favours configurations
        that preserve the full multiplicity of open directions without premature
        selection.
        Translated to the gauge connection, this principle selects configurations that
        minimise incompatibility — i.e. minimise $\|F_{\mu\nu}\|^2$.
        The action~\eqref{eq:ym-action} is the unique gauge-invariant, Lorentz-invariant,
        second-order functional of $A_\mu$ that achieves this.
      \end{proof}

      \begin{corollary}[Yang--Mills equations]
        Variation of~\eqref{eq:ym-action} with respect to $A_\mu$ yields
        \begin{equation}
          \label{eq:ym-equations}
          D_\mu F^{\mu\nu} \;=\; 0,
        \end{equation}
        the Yang--Mills equations in vacuum.
      \end{corollary}

      \begin{remark}[Status]
        Proposition~\ref{prop:ym-action} is a structural motivation, not a complete
        proof.
        The uniqueness claim requires showing that~\eqref{eq:ym-action} is the
        \emph{only} admissibility-compatible action at second order in $A_\mu$.
        This follows in the Born--Infeld limit from the results of~\cite{BornInfeld}
        (the BI action reduces to the Yang--Mills action in the weak-field regime),
        but a direct derivation from the axioms of~\cite{FoundationM} alone is
        an open step identified in Section~\ref{sec:open}.
      \end{remark}

  \section{Effective Spacetime Geometry}
    \label{sec:geometry}

    \subsection{Metric from the principal symbol}

      The operator-metric correspondence~\eqref{eq:metric-from-symbol} gives the
      effective metric as
      \begin{equation}
        g^{\mu\nu}(x) \;\propto\; A^{\mu\nu}(x).
      \end{equation}
      This metric is not postulated: it encodes how admissible modes propagate in
      the continuum limit of the projection.

    \subsection{Signature selection from admissibility}

      \begin{theorem}[Lorentzian signature from spectral admissibility,
        after~\cite{BornInfeld}]
        \label{thm:signature}
        In homogeneous and quasi-isotropic admissible regimes, the Born--Infeld
        admissibility constraint (A3--A4 of~\cite{FoundationM}) forces the effective
        metric to have Lorentzian signature $(-,+,+,+)$.
      \end{theorem}

      \begin{proof}
        The bounded Born--Infeld flux constraint excludes purely diffusive
        (parabolic) effective descriptions, which would imply instantaneous
        propagation of perturbations and contradict the finite capacity bound of A4.
        A Riemannian signature $(+,+,+,+)$ produces an elliptic principal symbol,
        which is also excluded.
        An ultra-hyperbolic signature (more than one time direction) fails to
        admit a globally well-posed initial value problem and is excluded by the
        spectral stability requirement.
        The unique remaining option compatible with (i) hyperbolicity of
        $\sigma_2(\mathcal{L}_\Pi)$, (ii) finite characteristic cone, (iii) the
        monotonic temporal ordering of Foundation M~§4.3, and (iv) the bounded
        admissibility constraint of A4, is the Lorentzian signature $(-,+,+,+)$.
        The detailed argument is established in~\cite{BornInfeld} Sections~5--6,
        which depend only on the bounded-flux constraint and spectral stability —
        both present in the Foundation M axioms — and not on any dynamical
        substrate.
      \end{proof}

    \subsection{Schwarzschild geometry from flux conservation}

      \begin{theorem}[Schwarzschild from bounded admissibility,
        after~\cite{BornInfeld}]
        \label{thm:schwarzschild}
        In the presence of a localised stationary obstruction to the admissibility
        flow, with spherical symmetry and no additional microscopic scales, the
        unique exterior stationary solution of the effective geometry induced by
        $\mathcal{L}_\Pi$ is the Schwarzschild metric
        \begin{equation}
          \label{eq:schwarzschild}
          ds^2
          \;=\; -\!\left(1 - \frac{r_s}{r}\right)dt^2
          + \left(1 - \frac{r_s}{r}\right)^{-1}dr^2
          + r^2 d\Omega^2,
        \end{equation}
        where $r_s$ is an integration constant characterising the strength of the
        obstruction.
      \end{theorem}

      \begin{proof}
        In the exterior region, the homogeneous equation for the effective
        operator~\eqref{eq:continuum-op} under spherical symmetry reduces to
        \begin{equation}
          \frac{1}{r^2}\frac{d}{dr}\!\left(r^2 A^r(r)\frac{d\Phi}{dr}\right) = 0,
        \end{equation}
        expressing conservation of admissible flux through spheres of radius $r$.
        The general solution is $\Phi(r) = \Phi_0 - C/r$, giving a $1/r$ decay profile.
        Via the operator-metric correspondence~\eqref{eq:metric-from-symbol}, this
        profile translates to the metric coefficients $g_{tt} = -f(r)$,
        $g_{rr} = f(r)^{-1}$ with $f(r) = 1 - r_s/r$.
        The parameter $r_s$ appears as the integration constant of the flux
        conservation equation and is identified operationally via its influence on
        admissible propagation cones.
        This derivation requires only: (i) the bounded-flux constraint (A3--A4),
        (ii) local injectivity of the projection in the exterior region,
        (iii) spherical symmetry and stationarity, (iv) absence of additional
        microscopic scales.
        No independent metric field equations or energy-momentum tensor
        are postulated.
        The detailed argument is established in~\cite{BornInfeld} Sections~6--8.
      \end{proof}

      \begin{remark}[Horizon as projection degeneracy]
        At $r = r_s$, the lapse function $f(r_s) = 0$.
        In the present framework, this corresponds to a degeneracy of the principal
        symbol $\sigma_2(\mathcal{L}_\Pi)$: the operator loses hyperbolicity and
        the continuum chart ceases to be projectable.
        The horizon is therefore not a physical singularity of the underlying
        admissible structure — it is the breakdown of the effective continuum
        description.
        This interpretation is consistent with the non-injectivity result of~\cite{ENI}:
        at the horizon, distinct microscopic configurations become maximally
        indistinguishable, maximising the projection entropy $S_{\Pi}$.
      \end{remark}

    \subsection{Einstein equations as consistency conditions}

      \begin{remark}[Einstein equations from coherence functional]
        The white paper~\cite{Cosmochrony} establishes, via the infrared expansion of
        the coherence functional (Appendix~A.10), that the Einstein--Hilbert action
        arises as the leading two-derivative infrared contribution.
        The Einstein equations $G_{\mu\nu} = 8\pi T_{\mu\nu}$ therefore emerge as
        consistency conditions of emergent geometry — not as microscopic laws of the
        admissible substrate.
        In the language of the present paper: they express the balance between the
        curvature of the projection-induced metric (encoding admissibility constraints)
        and the energy-momentum content (encoding localised obstructions to the
        admissibility flow).
      \end{remark}

  \section{Open Directions}
    \label{sec:open}

    \paragraph{Rigorous continuum limit.}
      The convergence $\Pi_n \to \mathcal{L}_\Pi$ for the Heisenberg Cayley graph
      $\mathrm{Cay}(G_q, S_q)$ in the $q \to \infty$ limit is established qualitatively
      in~\cite{FoundationM}.
      A fully rigorous proof at the level of the spectral O-series data would close
      this gap.

    \paragraph{Direct derivation of Yang--Mills uniqueness.}
      Proposition~\ref{prop:ym-action} gives a structural motivation for the
      Yang--Mills action.
      Its uniqueness — that $S_{\mathrm{YM}}$ is the \emph{only} admissibility-compatible
      second-order action — follows from the Born--Infeld limit~\cite{BornInfeld}
      but not yet from Foundation M axioms alone.

    \paragraph{SU(3) from admissibility.}
      The identification of $\mathrm{SU}(3)$ as a gauge group from the admissibility
      flow requires finding a co-admissibility pattern for an $\mathrm{SU}(3)$-type
      group in the O-series spectral data.
      The analogue of the $\lambda_{1/2} = \lambda_{3/2}$ result on $2I$ (established
      in~\cite{Q2}) must be found for a group with $\mathrm{SU}(3)$ structure.

    \paragraph{Coupling constants.}
      The present paper identifies which gauge groups and geometries are admissible.
      It does not determine the numerical values of coupling constants ($g_s$, $g_w$,
      $G_N$, $\alpha$).
      These are expected to emerge from the quantitative spectral data of the O-series
      via the transfer chain $c_\chi \to \delta_{\mathrm{pair}} \to \beta^* \approx 0.126$,
      but the explicit identification is open.

    \paragraph{Cosmological implications.}
      In homogeneous regimes, the admissibility flow gives a linear growth of cumulative
      projected capacity $I(n)$, underlying the emergent Hubble law and effective
      cosmological expansion~\cite{Cosmochrony}.
      The precise connection between the spectral exponent $\delta_{\mathrm{pair}}$
      and the observable cosmological parameters ($H_0$, $w$) is an open quantitative
      problem.

  \section*{Acknowledgements}

  The author acknowledges the use of large language models as a supportive tool
  for refining language, structure, and internal consistency during the development
  of this manuscript.
  All conceptual contributions, theoretical choices, and interpretations remain
  the sole responsibility of the author.

  \begin{thebibliography}{99}

    \bibitem{FoundationM}
    J.~Beau,
    ``Admissible Non-Injective Transitions as the Primitive of Physical Description,''
    Cosmochrony companion paper~M, Zenodo (2026).
    \doi{10.5281/zenodo.XXXXXXXX}.

    \bibitem{Q1}
    J.~Beau,
    ``Phase Coherence as an Admissibility Constraint:
    Structural Derivation from Born--Infeld Fibre Minimality,''
    Cosmochrony companion paper~M2, Zenodo (2026).
    \doi{10.5281/zenodo.XXXXXXXX}.

    \bibitem{Q2}
    J.~Beau,
    ``Quantum Structure beyond Spin-$\tfrac{1}{2}$:
    Co-Admissible Sectors and the Emergence of SU(2) as a Fixed Point,''
    Cosmochrony companion paper~M3, Zenodo (2026).
    \doi{10.5281/zenodo.XXXXXXXX}.

    \bibitem{ENI}
    J.~Beau,
    ``Non-Injectivity as a Structural Necessity of Genuine Emergence,''
    Cosmochrony companion paper~L, Zenodo (2026).
    \doi{10.5281/zenodo.19391746}.

    \bibitem{BornInfeld}
    J.~Beau,
    ``Bounded Relational Dynamics and the Structural Necessity of Born--Infeld Geometry,''
    Cosmochrony companion paper~C, Zenodo (2026).
    \doi{10.5281/zenodo.18407505}.

    \bibitem{Cosmochrony}
    J.~Beau,
    ``Cosmochrony: A Pre-Geometric Relational Framework for Emergent Physics,''
    white paper, Zenodo (2026).
    \doi{10.5281/zenodo.17957509}.

    \bibitem{O12}
    J.~Beau,
    ``Weil Representation Spectral Capacity on Heisenberg Cayley Graphs,''
    O12 spectral admissibility sub-programme paper, Zenodo (2026).
    \doi{10.5281/zenodo.19243611}.

    \bibitem{O23}
    J.~Beau,
    ``Three Stable Directions from Quaternionic Minimality,''
    O23 spectral admissibility sub-programme paper, Zenodo (2026).
    \doi{10.5281/zenodo.19375136}.

    \bibitem{O24}
    J.~Beau,
    ``Vertical Non-Injectivity and the Stability of the Observable Rank,''
    O24 spectral admissibility sub-programme paper, Zenodo (2026).
    \doi{10.5281/zenodo.19386581}.

  \end{thebibliography}

\end{document}
